\chapter*{General Introduction}
\addcontentsline{toc}{chapter}{General Introduction}

Security operations now take place under a form of time compression that did not exist a few years ago. Vulnerabilities are disclosed at high volume, proof-of-concept code appears quickly, cloud and SaaS identities have become privileged attack surfaces, and threat actors increasingly automate reconnaissance, credential theft, and social engineering. The limiting factor for many teams is no longer access to raw security data; it is the ability to correlate heterogeneous signals quickly enough to reach a defensible decision before conditions change again. Reports from CrowdStrike, Unit 42, Sophos, GitHub, Wiz, and CISA all point in the same direction: the operational window available to defenders has narrowed, while the number of relevant signals has expanded \cite{crowdstrike_gtr_2026,unit42_ir_2026,sophos_threat_2025,github_oss_trends_2026,wiz_ai_readiness_2026,cisa_kev}.

The platform studied in this report was engineered for that environment. It is a self-hosted Threat Intelligence Platform (TIP) designed for a small security organization that needs one working system rather than a collection of disconnected products. The implementation combines fifteen FastAPI services, nine shared Python packages, a Next.js frontend with a backend-for-frontend proxy, PostgreSQL behind PgBouncer, Redis, a LiteLLM gateway, and a scheduler-driven execution model. Its purpose is not to collect threat data for its own sake. Its purpose is to reduce the distance between raw intelligence and operational judgment: which vulnerabilities matter to this organization, which threat actors are plausible, which indicators deserve attention now, and what should an analyst or manager do next.

Several architectural decisions follow directly from that operational objective. Intelligence ingestion is separated from AI-driven synthesis so that data collection survives model outages and quota exhaustion. Services own their schemas and communicate through HTTP rather than shared table access, which preserves modularity even within a single-host deployment. Secrets are centralized, outbound calls are wrapped in a resilience layer, and the latency-critical indicator lookup path is cached aggressively. The implementation therefore reflects a specific engineering stance: availability, traceability, and operability matter at least as much as algorithmic sophistication.

This report rebuilds the project narrative from the implementation upward. It does not treat the source code as an appendix to a design story written in advance. Instead, the architecture, workflows, trust boundaries, deployment model, testing posture, and operational tradeoffs are derived from the repository and its infrastructure. The first chapter situates the work in the modern threat landscape and explains why partially automated, AI-assisted security operations have become operationally necessary. The second establishes the theoretical foundations required to understand the project properly. The third positions the work against existing solutions. The fourth explains the methodology and engineering process through which the system was analysed and built. The fifth translates the problem context into requirements and design constraints. The sixth explains the architecture of the delivered system. The seventh clarifies the infrastructure stack and platform technologies in detail. The eighth introduces the service catalogue and capability map. The ninth deepens that view with a service-by-service architectural reading. The tenth explains the data model, API design, and frontend workflows. The eleventh examines the implementation, runtime behaviour, deployment model, and operational mechanics. The final chapter addresses verification, testing, limitations, and the most credible next steps toward a production-hardened platform.

This new structure is intentional. A short engineering report can jump quickly from context to implementation because it assumes the reader already shares the design vocabulary of the author. A thesis cannot make that assumption. It must explain the conceptual foundations that justify the chosen abstractions, the tradeoffs that differentiate the solution from neighbouring designs, and the limitations that prevent inflated claims. For that reason, the report explicitly introduces theoretical foundations between the threat chapter and the solution chapters. This enlarges the document, but more importantly it makes the argument defensible: the reader can evaluate not only what was built, but why those design choices are intellectually coherent.

Methodologically, the report adopts a code-grounded reverse-engineering approach. The repository, infrastructure definitions, runtime configuration, rendered interface evidence, and internal engineering notes are treated as primary artefacts. Public threat reports and standards are used as external context, not as substitutes for implementation evidence. This matters because the most common weakness in student security reports is narrative drift: the text slowly begins to describe an idealized platform rather than the one that actually exists. The present report is designed specifically to avoid that failure mode.
